\subsection{Sturm-Liouville Boundary-Value Problems}\label{sec:sturm-liouville-boundary-value-problems}

We now study a class of second-order ODE boundary-value problems which arise from separation of variables. A \textbf{Sturm-Liouville problem} (or S-L problem) involves the ordinary differential equation
\begin{equation}\label{eq:sturm-liouville-inhomogeneous-differential-equation}
-\frac{\mathrm{d}}{\mathrm{d}x}\left(p(x)\frac{\mathrm{d}u}{\mathrm{d}x}(x)\right)+q(x)u(x)-\lambda w(x)u(x)=f(x)
\end{equation}
on the interval $(a,b)$ and appropriate boundary conditions which we describe below. We assume the following:\begin{enumerate}\item The functions $p$, $p\prime$, $q$ and $w$ are real-valued and continuous on the open interval $(a,b)$.\item The functions $p$ and $w$ are positive on $(a,b)$.\end{enumerate}We say the S-L problem is \textbf{regular} if both $a$ and $b$ are finite and assumptions 1 and 2 hold on the closed interval $[a,b]$. If not, we say the problem is \textbf{singular}.

We formally define the differential operator
\begin{equation}\label{eq:sturm-liouville-weighted-differential-operator}
Lu:=\dfrac{1}{w}[-(pu')'+qu],
\end{equation}
and we note that \eqref{eq:sturm-liouville-inhomogeneous-differential-equation} can now be written in the form
\begin{equation}\label{eq:sturm-liouville-weighted-operator-equation}
Lu-\lambda u=\dfrac{f}{w}.
\end{equation}
We intend to use the theory just developed above to analyze this as an eigenvalue problem for an operator from the weighted space $L_w^2(a,b)$ to itself. However, since the analysis of singular problems emphasizes methods other than those we have described we will discuss only regular problems.

We use the weighted space $L_w^2(a,b)$, but in regular problems this is really nothing more than a notational convenience since
\begin{equation}\label{eq:sturm-liouville-weighted-norm-equivalence}
\min_{x\in[a,b]}w(x)\|u\|_{L^2(a,b)}^2\leq\|u\|_{L_w^2(a,b)}^2\leq\max_{x\in[a,b]}w(x)\|u\|_{L^2(a,b)}^2,
\end{equation}
so that the $L^2$ and $L_w^2$ norms are equivalent. We will use this to define domains for the operator $L$.

We will examine the most common type of boundary conditions for S-L problems encountered in applications, namely, those of \textbf{unmixed type}. We require
\begin{equation}\label{eq:sturm-liouville-left-unmixed-boundary-condition}
\cos\alpha\,u(a)-\sin\alpha\,u'(a)=0,
\end{equation}

\begin{equation}\label{eq:sturm-liouville-right-unmixed-boundary-condition}
\cos\beta\,u(b)+\sin\beta\,u'(b)=0.
\end{equation}
We now define the domain
\begin{equation}\label{eq:sturm-liouville-unmixed-operator-domain}
\mathcal{D}(L):=\{u\in H^2(a,b)\mid\eqref{eq:sturm-liouville-left-unmixed-boundary-condition}\text{ and }\eqref{eq:sturm-liouville-right-unmixed-boundary-condition}\text{ are satisfied}\}.
\end{equation}
We now prove the following theorem.

\begin{Theorem}
\label{thm:regular-sturm-liouville-first-spectral-properties}
Let $(\mathcal{D}(L),L)$ be defined by \eqref{eq:sturm-liouville-weighted-differential-operator} and \eqref{eq:sturm-liouville-unmixed-operator-domain}. The following hold:\begin{enumerate}\item The eigenvalues of $(\mathcal{D}(L),L)$ are real.\item The eigenvalues of $(\mathcal{D}(L),L)$ are bounded below by a constant $\lambda_G\in\mathbb{R}$.\item Eigenfunctions corresponding to distinct eigenvalues are mutually orthogonal in $L_w^2(a,b)$.\item Each eigenvalue has multiplicity one.\end{enumerate}
\end{Theorem}

\begingroup
\postdisplaypenalty=10000
\renewcommand{\endProof}{\par\nobreak\hfill$\blacksquare$\par\vspace{0.1cm}}
\begin{Proof}
To begin, we integrate by parts to prove Lagrange\textquotesingle s identity; i.e., that for every $u$ and $v$ in $H^2(a,b)$ we have
\begin{equation}\label{eq:sturm-liouville-lagrange-identity}
(Lu,v)_w-(u,Lv)_w=\left.p(x)[u(x)v'(x)-u'(x)v(x)]\right|_a^b.
\end{equation}
Thus, if $u$ and $v$ are in $\mathcal{D}(L)$ we can use the boundary conditions \eqref{eq:sturm-liouville-left-unmixed-boundary-condition} and \eqref{eq:sturm-liouville-right-unmixed-boundary-condition} to get
\begin{equation}\label{eq:sturm-liouville-domain-symmetry}
(Lu,v)_w=(u,Lv)_w,
\end{equation}
proving that $(\mathcal{D}(L),L)$ is symmetric. Hence, Theorem~\ref{thm:symmetric-operator-spectral-properties} immediately gives us parts 1 and 3.

To prove part 2 we prove an energy estimate of the form
\begin{equation}\label{eq:sturm-liouville-lower-energy-bound}
(Lu,u)_w\geq\lambda_G\|u\|_{L_w^2(a,b)}^2
\end{equation}
for all $u\in\mathcal{D}(L)$. (This is an analogue of Gårding\textquotesingle s inequality in elliptic PDEs (cf. Section~\ref{subsec:elliptic-garding-inequality}), hence the notation $\lambda_G$.) We prove this only in the case $\tan\alpha,\tan\beta\in[0,\infty)$ and leave the proof of other cases to the reader (Problem~\ref{ex:sturm-liouville-general-energy-bound}). For any $u\in\mathcal{D}(L)$ we have\[\begin{aligned}
(Lu,u)&=\int_a^b-(pu')'u+q|u|^2\,\mathrm{d}x\\
&=\int_a^bp|u'|^2+q|u|^2\,\mathrm{d}x+p(a)u'(a)u(a)-p(b)u'(b)u(b)\\
&=\int_a^bp|u'|^2+q|u|^2\,\mathrm{d}x\\
&\quad+p(a)\tan\alpha|u'(a)|^2+p(b)\tan\beta|u'(b)|^2\\
&\geq\int_a^bq(x)|u(x)|^2\,\mathrm{d}x\\
&\geq\dfrac{\min_{x\in[a,b]}q(x)}{\max_{x\in[a,b]}w(x)}\|u\|_{L_w^2(a,b)}^2.
\end{aligned}\]To get part 2 we simply observe that for any eigenvalue $\lambda$ we have
\begin{equation}\label{eq:sturm-liouville-eigenfunction-energy-identity}
(Lu,u)_w=(\lambda wu,u)=\lambda(u,u)_w.
\end{equation}
Hence $\lambda\geq\lambda_G$.

Part 4 follows immediately from the uniqueness theorem for initial-value problems for ODEs, which implies that either of the boundary conditions \eqref{eq:sturm-liouville-left-unmixed-boundary-condition} or \eqref{eq:sturm-liouville-right-unmixed-boundary-condition} determines a solution of the homogeneous ODE $Lu=\lambda u$ up to a multiplicative constant.
\end{Proof}
\endgroup
